\section{CNN 训练实用技巧}

\subsection{学习率调度}

\begin{figure}[H]
\centering
\includegraphics[width=0.95\textwidth]{slides-images/slide-039.jpg}
\caption{常见的学习率调度策略：阶梯衰减、余弦退火、线性预热\protect\footnotemark}
\end{figure}
\footnotetext{Slides 第40页。}

学习率是训练中最关键的超参数之一。常见的调度策略：

\begin{itemize}
    \item \textbf{阶梯衰减}（Step Decay）：每隔固定 epoch 将学习率乘以衰减系数（如 0.1）
    \item \textbf{余弦退火}（Cosine Annealing）：学习率按余弦曲线从初始值平滑衰减到接近零
    \item \textbf{线性预热}（Linear Warmup）：训练初期从很小的学习率线性增加到目标学习率
    \item \textbf{预热 + 余弦退火}：先线性预热，再余弦衰减——这是现代训练中最常用的策略
\end{itemize}

\begin{importantbox}{学习率调度的直觉}
\begin{itemize}
    \item \textbf{初期需要预热}：刚开始时，网络权重随机，梯度方向不稳定。大学习率可能导致训练不稳定。
    \item \textbf{中期需要较大学习率}：快速探索参数空间，找到好的区域。
    \item \textbf{后期需要小学习率}：在好的区域内精细调整，收敛到更优的解。
\end{itemize}
\end{importantbox}

\subsection{数据增强}

\begin{figure}[H]
\centering
\includegraphics[width=0.95\textwidth]{slides-images/slide-041.jpg}
\caption{常见的数据增强技术：随机裁剪、水平翻转、颜色抖动等\protect\footnotemark}
\end{figure}
\footnotetext{Slides 第42页。}

数据增强通过对训练图像施加随机变换来人工扩大数据集，是最有效的正则化技术之一：

\begin{itemize}
    \item \textbf{随机裁剪}（Random Crop）：在略大的图像中随机裁出目标大小的区域
    \item \textbf{水平翻转}（Horizontal Flip）：50\% 概率左右镜像翻转
    \item \textbf{颜色抖动}（Color Jitter）：随机调整亮度、对比度、饱和度
    \item \textbf{随机擦除}（Random Erasing）：随机遮挡图像的一部分
    \item \textbf{Mixup / CutMix}：将两张图像混合或拼接
\end{itemize}

\begin{knowledgebox}{数据增强的核心原则}
数据增强的变换应该满足：
\begin{itemize}
    \item \textbf{标签不变性}：变换后的图像应该仍然属于同一类别（水平翻转的猫还是猫）
    \item \textbf{合理性}：变换应该生成\textbf{可能出现在现实中}的图像
    \item \textbf{训练时使用，测试时不用}（或使用确定性版本，如中心裁剪）
\end{itemize}
\end{knowledgebox}

\subsection{正则化方法汇总}

\begin{figure}[H]
\centering
\includegraphics[width=0.95\textwidth]{slides-images/slide-043.jpg}
\caption{正则化方法汇总：L2 权重衰减、Dropout、数据增强等\protect\footnotemark}
\end{figure}
\footnotetext{Slides 第44页。}

防止过拟合的主要手段：
\begin{itemize}
    \item \textbf{L2 正则化 / 权重衰减}：损失函数中加入 $\lambda \|W\|^2$ 惩罚项
    \item \textbf{Dropout}：训练时随机关闭神经元
    \item \textbf{数据增强}：增加训练数据的多样性
    \item \textbf{提前停止}（Early Stopping）：在验证集性能开始下降时停止训练
    \item \textbf{Batch Normalization}：也有一定的正则化效果
\end{itemize}

\subsection{迁移学习}

\begin{figure}[H]
\centering
\includegraphics[width=0.95\textwidth]{slides-images/slide-045.jpg}
\caption{迁移学习：在大数据集上预训练，然后在小数据集上微调\protect\footnotemark}
\end{figure}
\footnotetext{Slides 第46页。}

当目标任务的数据量有限时，迁移学习是最有效的策略：

\begin{enumerate}
    \item 在大规模数据集（如 ImageNet）上\textbf{预训练}一个 CNN
    \item 将预训练模型的卷积层作为\textbf{特征提取器}
    \item 替换最后的全连接层（匹配目标任务的类别数）
    \item 用目标数据集\textbf{微调}（fine-tune）——可以只训练最后几层，也可以全网络微调（通常用较小的学习率）
\end{enumerate}

\begin{importantbox}{迁移学习的经验法则}
\begin{itemize}
    \item 数据很少 + 与预训练域相似：只训练最后一层
    \item 数据较多 + 与预训练域相似：微调全网络（小学习率）
    \item 数据较多 + 与预训练域不同：微调全网络（较大学习率）
    \item 数据很少 + 与预训练域不同：最困难的情况，考虑使用更多数据增强
\end{itemize}
\end{importantbox}

\subsection{本章小结}

训练一个高性能的 CNN 不仅需要好的架构设计，还需要一系列工程技巧的配合：Kaiming 初始化确保训练稳定启动，学习率调度控制训练节奏，数据增强增强泛化能力，正则化防止过拟合，迁移学习利用已有知识加速新任务的学习。

%% ========== 总结与延伸 ==========

